%% chapters/chapter01_introduction.tex
%% Capítulo 1: Introducción a los Sistemas de Control
%% Apunte "Control Automático" — UTN IEE
%% Codificación: UTF-8

\chapter{Introducción a los Sistemas de Control}
\label{ch:introduccion}

\begin{flushright}
  \itshape
  ``El propósito del control automático es hacer que el comportamiento\\
  de un sistema se ajuste a un objetivo deseado, sin intervención humana\\
  continua, incluso frente a perturbaciones imprevistas.''\\[2pt]
  \normalfont\small --- Katsuhiko Ogata, \textit{Ingeniería de Control Moderna}
\end{flushright}

\bigskip

El control automático está en todas partes: en el termostato que regula la
temperatura de tu hogar, en el regulador de velocidad del automóvil, en
los convertidores de frecuencia que controlan motores industriales, en los
sistemas que estabilizan la tensión de la red eléctrica.
En este capítulo se introducen los conceptos fundamentales que servirán
de base para todo el curso. Al finalizar, el estudiante será capaz de:
\begin{enumerate}
  \item Identificar los componentes de un sistema de control en cualquier
        aplicación práctica.
  \item Distinguir entre control de lazo abierto y lazo cerrado, y comprender
        las ventajas de la realimentación.
  \item Describir cualitativamente las especificaciones de desempeño de un
        sistema de control.
  \item Implementar simulaciones básicas en MATLAB para analizar la respuesta
        de sistemas simples.
\end{enumerate}

%======================================================================
\section{¿Qué es el control automático?}
\label{sec:intro-definicion}
%======================================================================

Imaginá que tenés que mantener la velocidad de un motor eléctrico exactamente
en 1500 RPM mientras la carga mecánica que impulsa varía de manera impredecible.
Si no hacés nada, el motor se frenará o se acelerará con cada cambio de carga.
Podrías encargarte de ajustar manualmente la tensión de alimentación en cada
instante, pero eso es humanamente imposible con la rapidez y precisión que requiere
la aplicación. La solución es un \textbf{sistema de control automático}: un
mecanismo que, sin intervención humana permanente, regula la velocidad del motor
para mantenerla en el valor deseado pese a las perturbaciones.

\begin{definicion}{Sistema de control automático}
  Un \textbf{sistema de control automático} es un conjunto de componentes físicos
  y/o lógicos que, mediante la observación de la salida de un proceso y la
  comparación con un valor deseado, genera acciones correctivas para lograr que
  dicha salida siga el valor deseado con la precisión, velocidad y estabilidad
  especificadas, reduciendo o eliminando el efecto de perturbaciones externas.
\end{definicion}

La palabra clave en esa definición es \textit{observación}: el sistema de control
\textbf{mide} lo que está sucediendo, lo \textbf{compara} con lo que debería suceder
y \textbf{actúa} en consecuencia. Este ciclo ---medir, comparar, actuar--- es la
esencia del control automático y se denomina \textbf{realimentación} o
\textbf{retroalimentación} (del inglés \textit{feedback}).

\subsection{¿Por qué ``automático''?}

El adjetivo \textit{automático} distingue estos sistemas de los que requieren
ajuste humano constante. En un control manual, un operario observa un medidor,
juzga si la variable está fuera del rango deseado y opera algún actuador para
corregirla. Este esquema:
\begin{itemize}
  \item Depende de la atención y el tiempo de reacción humanos.
  \item No puede seguir perturbaciones más rápidas que unos pocos segundos.
  \item No puede controlar simultáneamente decenas de variables en un proceso
        industrial complejo.
\end{itemize}
Un sistema de control automático elimina estas limitaciones. Puede reaccionar
en microsegundos, operar sin fatiga las 24 horas del día y controlar
simultáneamente miles de variables en una planta industrial.

\subsection{Diferencia entre regulación y seguimiento}

Existen dos objetivos distintos que puede tener un sistema de control:

\begin{description}
  \item[Regulación (control de regulación):] mantener la variable controlada
    constante frente a perturbaciones. Ejemplo: mantener la tensión de un
    generador en 380 V pese a variaciones de carga.
  \item[Seguimiento (control servo):] hacer que la variable controlada
    siga una referencia que cambia con el tiempo. Ejemplo: un brazo robótico
    que sigue una trayectoria predefinida.
\end{description}

Los mismos principios y herramientas matemáticas se aplican a ambos casos,
aunque el énfasis en el diseño puede variar.

%======================================================================
\section{Perspectiva histórica}
\label{sec:intro-historia}
%======================================================================

Comprender cómo evolucionó el control automático ayuda a apreciar por qué
se desarrollaron ciertas herramientas matemáticas y cuál es su relevancia práctica.

\textbf{Siglo XVIII — El regulador centrífugo de Watt (1788).}
James Watt incorporó un regulador centrífugo a su máquina de vapor para
mantener la velocidad constante. Dos bolas giratorias suben o bajan por
efecto centrífugo a medida que la velocidad varía, y ese movimiento acciona
una válvula de vapor. Es el primer sistema de control realimentado de uso
industrial masivo.

\textbf{Siglo XIX — Análisis matemático.}
El matemático James Clerk Maxwell analizó en 1868 la estabilidad del regulador
de Watt, publicando el primer tratamiento matemático de la estabilidad de
sistemas de control. Eduard Routh (1877) generalizó este análisis con el
criterio algebraico que lleva su nombre \verCap{7}.

\textbf{Siglo XX temprano — Control PID y automatización industrial.}
Nicolas Minorsky formuló en 1922 el controlador proporcional-integral-derivativo
(PID) en el contexto del control de dirección de barcos. El PID sigue siendo,
hasta hoy, el controlador más usado en la industria.

\textbf{1930s--1940s — Métodos de dominio frecuencial.}
Harry Nyquist (1932) desarrolló su criterio de estabilidad basado en la
respuesta en frecuencia. Hendrik Bode (1940) sistematizó el análisis de
frecuencia con sus diagramas logarítmicos, imprescindibles para el diseño
de compensadores \verCap{10}.

\textbf{1960s — Control moderno y espacio de estados.}
Rudolf Kálmán introdujo el enfoque de variables de estado y el filtro de
Kálmán, abriendo la era del control moderno. Estos métodos permiten analizar
sistemas de múltiples entradas y múltiples salidas (MIMO) \verCap{14}.

\textbf{1980s-presente — Control digital y computacional.}
El control se implementa hoy casi exclusivamente en procesadores digitales
(DSPs, PLCs, microcontroladores). MATLAB/Simulink se convirtió en la
herramienta estándar para análisis, diseño y simulación.

\begin{observacion}[Relevancia para Ingeniería Eléctrica]
  La mayoría de los hitos históricos del control automático surgieron de
  problemas de ingeniería eléctrica y mecánica: motores, máquinas de vapor,
  comunicaciones. Hoy, los sistemas eléctricos de potencia, la electrónica de
  potencia y las energías renovables son los principales impulsores de nuevas
  técnicas de control.
\end{observacion}

%======================================================================
\section{El sistema de control en la ingeniería eléctrica}
\label{sec:intro-aplicaciones-ee}
%======================================================================

La ingeniería en energía eléctrica es uno de los campos donde el control
automático tiene mayor impacto. A continuación se presentan las áreas de
aplicación más relevantes, que se retomarán como ejemplos a lo largo del curso.

\subsection{Control de accionamientos eléctricos}

Un \textbf{accionamiento eléctrico} (o \textit{drive}) es el conjunto formado
por un motor eléctrico y el convertidor electrónico que lo alimenta.
El control automático es fundamental en:

\begin{itemize}
  \item \textbf{Control de velocidad:} un variador de frecuencia (inversor)
    regula la velocidad de un motor de inducción ajustando la frecuencia
    y tensión de alimentación. Sin control realimentado, la velocidad variaría
    con la carga.
  \item \textbf{Control de par:} en tracción eléctrica (trenes, autos eléctricos),
    el control de par permite aceleraciones suaves y respuesta rápida.
  \item \textbf{Control de posición (servo):} en brazos robóticos y máquinas
    CNC, el motor debe posicionarse con precisión de micrómetros.
\end{itemize}

\begin{aplicacion}[Variador de frecuencia en bomba industrial]
  Una bomba centrífuga alimenta agua a un edificio. Sin control, la bomba
  funciona a velocidad constante y la presión se regula con una válvula de
  estrangulamiento (ineficiente). Con un variador de frecuencia controlado
  por un sensor de presión diferencial, la velocidad de la bomba se ajusta
  continuamente para mantener la presión de consigna. El ahorro energético
  típico es del 30--50\%.
\end{aplicacion}

\subsection{Control en sistemas de potencia}

Los sistemas eléctricos de potencia ---generación, transmisión y distribución---
dependen de múltiples lazos de control para su operación estable:

\begin{itemize}
  \item \textbf{Regulador automático de tensión (AVR):} controla la corriente
    de excitación de un generador sincrónico para mantener la tensión terminal
    constante.
  \item \textbf{Regulador de velocidad (gobernador):} controla la apertura de
    las válvulas de entrada de vapor (turbina) o agua (hidráulica) para mantener
    la frecuencia de la red en \SI{50}{\hertz}.
  \item \textbf{Estabilizadores de sistema de potencia (PSS):} añaden
    amortiguamiento a las oscilaciones electromecánicas de la red.
  \item \textbf{Control de reactivos (FACTS):} dispositivos de electrónica de
    potencia que controlan el flujo de potencia reactiva para mejorar la
    estabilidad de tensión.
\end{itemize}

\subsection{Electrónica de potencia}

Los convertidores estáticos (rectificadores, inversores, choppers) requieren
lazos de control muy rápidos (tiempos de respuesta de microsegundos) para:

\begin{itemize}
  \item Mantener la tensión de salida constante frente a variaciones de carga.
  \item Regular la corriente para proteger los semiconductores.
  \item Inyectar corriente senoidal a la red (inversores fotovoltaicos).
  \item Controlar el flujo de potencia en enlace de corriente continua de alta
    tensión (HVDC).
\end{itemize}

\subsection{Energías renovables}

Los sistemas de energía renovable presentan desafíos particulares de control:

\begin{itemize}
  \item \textbf{Turbinas eólicas:} control del ángulo de pala (\textit{pitch})
    para extraer la máxima potencia del viento (seguimiento del punto de máxima
    potencia, MPPT) y limitar la potencia ante vientos fuertes.
  \item \textbf{Sistemas fotovoltaicos:} control MPPT mediante algoritmos
    de perturbación y observación que ajustan el punto de operación.
  \item \textbf{Almacenamiento en baterías:} control de carga/descarga para
    optimizar la vida útil y gestionar la inyección a la red.
\end{itemize}

\begin{nota}[Transversalidad del control automático]
  Todas las aplicaciones mencionadas comparten los mismos conceptos
  matemáticos: función de transferencia, estabilidad, respuesta en frecuencia,
  diseño de controladores. Las herramientas que se desarrollan en este curso
  son directamente aplicables a cada uno de estos campos.
\end{nota}

%======================================================================
\section{Componentes fundamentales de un sistema de control}
\label{sec:intro-componentes}
%======================================================================

Para analizar y diseñar sistemas de control, es indispensable identificar
correctamente sus componentes. Todo sistema de control (con realimentación)
contiene los siguientes elementos:

\subsection{Planta o proceso}

La \textbf{planta} (también llamada \textbf{proceso} en el contexto industrial)
es el objeto físico cuya variable de salida se quiere controlar. Es lo que
el sistema debe gobernar.

\textit{Ejemplos:} un motor eléctrico (salida: velocidad de rotación), un horno
(salida: temperatura), un tanque (salida: nivel de líquido), un generador
(salida: tensión terminal).

La planta tiene una dinámica propia ---descrita por ecuaciones diferenciales
o, en el dominio de Laplace, por una función de transferencia--- que determina
cómo responde a las señales de entrada.

\subsection{Actuador}

El \textbf{actuador} es el dispositivo que ejecuta la acción de control sobre
la planta. Convierte la señal de control (generalmente eléctrica, de baja
potencia) en una señal física de alta potencia capaz de modificar el
comportamiento de la planta.

\textit{Ejemplos:}
\begin{itemize}
  \item Variador de frecuencia (accionamiento eléctrico).
  \item Válvula motorizada (sistemas de fluidos o vapor).
  \item Transistor IGBT o MOSFET en un convertidor (electrónica de potencia).
  \item Servo motor con engranaje (robótica).
\end{itemize}

\subsection{Sensor o transductor}

El \textbf{sensor} mide la variable de salida de la planta y la convierte en
una señal eléctrica que puede compararse con la referencia.

\textit{Ejemplos:}
\begin{itemize}
  \item Tacómetro o encoder incremental (velocidad angular).
  \item Termopar o RTD (temperatura).
  \item Transformador de medida de tensión (tensión eléctrica).
  \item Transformador de corriente (corriente eléctrica).
  \item Resolver o encoder absoluto (posición angular).
\end{itemize}

La precisión del sensor determina un límite superior a la precisión del
sistema de control: \textit{no se puede controlar mejor que lo que se puede medir}.

\subsection{Controlador}

El \textbf{controlador} es el ``cerebro'' del sistema. Recibe la señal de error
(diferencia entre la referencia y la medición), la procesa matemáticamente y
genera la señal de control que actúa sobre el actuador.

Hoy en día, la mayor parte de los controladores son digitales: PLCs, DSPs,
microcontroladores o computadoras industriales que ejecutan algoritmos de
control en tiempo real.

Los tipos de controladores que se estudian en este curso incluyen:
controladores proporcionales (P), proporcionales-integrales (PI),
proporcionales-derivativos (PD), PID, compensadores en adelanto y en
atraso, y reguladores de estado.

\subsection{Comparador o elemento de suma}

El \textbf{comparador} (también llamado \textbf{sumador} o \textbf{nodo de
error}) calcula la diferencia entre la señal de referencia $r(t)$ y la señal
medida de salida $b(t)$:
\begin{equation}
  e(t) = r(t) - b(t)
  \label{eq:error}
\end{equation}
Esta diferencia $e(t)$ se llama \textbf{señal de error}. Es la entrada al
controlador. Si $e(t) = 0$, la salida coincide exactamente con la referencia
y el controlador no necesita generar ninguna corrección.

\subsection{Señales del sistema}

\begin{definicion}{Nomenclatura de señales}
  \begin{description}[leftmargin=3.5em, labelwidth=3em]
    \item[$r(t)$] \textbf{Señal de referencia} (o \textit{setpoint}, \textit{consigna}).
      Valor deseado de la variable controlada.
    \item[$e(t)$] \textbf{Señal de error}. $e(t) = r(t) - b(t)$.
    \item[$u(t)$] \textbf{Señal de control} (acción de control). Salida del
      controlador, entrada al actuador/planta.
    \item[$y(t)$] \textbf{Salida del sistema} (o variable controlada).
    \item[$b(t)$] \textbf{Señal de realimentación}. Medición de la salida
      proporcionada por el sensor: $b(t) = H \cdot y(t)$ donde $H$ es la
      ganancia del sensor.
    \item[$d(t)$] \textbf{Perturbación} (o disturbio, \textit{disturbance}).
      Señal no deseada que afecta la salida. Puede actuar sobre la planta
      o sobre el actuador.
    \item[$n(t)$] \textbf{Ruido de medición}. Señal que contamina la
      medición del sensor.
  \end{description}
\end{definicion}

La \cref{fig:componentes-sistema} muestra el diagrama de bloques con todos
estos componentes y señales.

% [INSERTAR FIGURA: Diagrama de bloques completo de un sistema de control
%  con realimentación negativa. Mostrar todos los componentes: referencia r(t),
%  sumador, controlador Gc(s), actuador Ga(s), planta Gp(s), sensor H(s).
%  Mostrar las señales e(t), u(t), y(t), b(t), perturbación d(t) y ruido n(t).
%  El alumno debe observar el flujo de señales y el lazo de realimentación.
%  Usar colores para distinguir: lazo directo (azul), lazo de realimentación (rojo).]
\begin{figure}[htbp]
  \centering
  \figplaceholder{Diagrama de bloques completo de un sistema de control realimentado.
  Componentes: referencia, sumador, controlador, actuador, planta, sensor.
  Señales: $r(t)$, $e(t)$, $u(t)$, $y(t)$, $b(t)$, $d(t)$, $n(t)$.}
  \caption{Diagrama de bloques completo de un sistema de control realimentado.}
  \label{fig:componentes-sistema}
\end{figure}

%======================================================================
\section{Sistema de control de lazo abierto}
\label{sec:intro-lazo-abierto}
%======================================================================

\begin{definicion}{Sistema de lazo abierto}
  Un \textbf{sistema de control de lazo abierto} es aquel en el que la acción
  de control es independiente de la salida del proceso. No existe medición
  de la salida ni comparación con la referencia: la señal de control se genera
  únicamente a partir de la referencia de entrada.
\end{definicion}

\subsection{Estructura de un sistema de lazo abierto}

En un sistema de lazo abierto, la estructura es simplemente:
\begin{equation}
  r(t) \;\longrightarrow\; \boxed{\text{Controlador}} \;\longrightarrow\;
  u(t) \;\longrightarrow\; \boxed{\text{Planta}} \;\longrightarrow\; y(t)
\end{equation}

La salida $y(t)$ no ``vuelve'' al controlador. El controlador actúa según un
programa prefijado, sin saber qué está pasando realmente con la planta.

\begin{figure}[htbp]
  \centering
  \begin{tikzpicture}[auto, node distance=2.8cm, thick]
    \node[input]  (ref)  {};
    \node[block,  right of=ref, node distance=2.5cm] (ctrl) {Controlador\\$G_c(s)$};
    \node[block,  right of=ctrl, node distance=3.8cm] (planta) {Planta\\$G_p(s)$};
    \node[output, right of=planta, node distance=2.5cm] (out) {};

    \draw[->] (ref)   -- node[above] {$R(s)$} (ctrl);
    \draw[->] (ctrl)  -- node[above] {$U(s)$} (planta);
    \draw[->] (planta)-- node[above] {$Y(s)$} (out);
  \end{tikzpicture}
  \caption{Diagrama de bloques de un sistema de control de \textbf{lazo abierto}.
           No hay medición de la salida ni realimentación.}
  \label{fig:lazo-abierto}
\end{figure}

\subsection{Ejemplo: temporizador de microondas}

Un horno microondas controlado por temporizador es un ejemplo clásico de
lazo abierto. El usuario fija el tiempo de calentamiento (referencia).
El horno opera durante ese tiempo y se detiene. No mide la temperatura
final de la comida ni ajusta la potencia si el alimento tardó más o menos
de lo esperado. El resultado depende enteramente de que el modelo inicial
(``tanta potencia durante tanto tiempo produce tal temperatura'') sea correcto.

\subsection{Ventajas del lazo abierto}

\begin{itemize}
  \item \textbf{Simplicidad:} no requiere sensor de salida, comparador ni
    algoritmo de realimentación. Es más barato y fácil de implementar.
  \item \textbf{Estabilidad garantizada:} si la planta es estable por sí
    misma, el lazo abierto también lo es. No existe el riesgo de oscilaciones
    inducidas por la realimentación.
  \item \textbf{Ausencia de ruido de realimentación:} no introduce el ruido
    del sensor al lazo de control.
\end{itemize}

\subsection{Desventajas del lazo abierto}

\begin{itemize}
  \item \textbf{Sensibilidad a perturbaciones:} si actúa una perturbación
    sobre la planta, la salida se desvía del valor deseado y nada la corrige.
  \item \textbf{Sensibilidad a variaciones de parámetros:} si los parámetros
    de la planta cambian (temperatura, envejecimiento, carga), la salida
    no coincidirá con la deseada.
  \item \textbf{Error de estado estacionario:} en general, la salida no
    coincide exactamente con la referencia salvo que el modelo de la planta
    sea perfecto (lo cual nunca ocurre en la práctica).
\end{itemize}

%======================================================================
\section{Sistema de control de lazo cerrado con realimentación}
\label{sec:intro-lazo-cerrado}
%======================================================================

\begin{definicion}{Sistema de lazo cerrado}
  Un \textbf{sistema de control de lazo cerrado} (o \textbf{sistema realimentado})
  es aquel en el que la acción de control depende de la salida del proceso.
  La salida se mide continuamente, se compara con la referencia y la diferencia
  (error) se usa para generar la acción correctiva.
\end{definicion}

El diagrama de bloques estándar del lazo cerrado con realimentación negativa
es el siguiente:

\begin{figure}[htbp]
  \centering
  \begin{tikzpicture}[auto, node distance=2.8cm, thick]
    %% Nodos
    \node[input]   (ref)  {};
    \node[sumnode, right of=ref, node distance=1.8cm] (suma) {};
    \node[block,   right of=suma, node distance=3.2cm] (ctrl)   {Controlador\\$G_c(s)$};
    \node[block,   right of=ctrl, node distance=3.8cm] (planta) {Planta\\$G_p(s)$};
    \node[output,  right of=planta, node distance=2.5cm] (out) {};
    \node[block,   below of=planta, node distance=1.8cm] (sensor) {Sensor\\$H(s)$};
    \node[junction, below of=out, node distance=1.8cm] (junc) {};

    %% Flechas del lazo directo
    \draw[->] (ref)    -- node[above, pos=0.5] {$R(s)$} (suma);
    \draw[->] (suma)   -- node[above] {$E(s)$} (ctrl);
    \draw[->] (ctrl)   -- node[above] {$U(s)$} (planta);
    \draw[->] (planta) -- node[above, name=yout, pos=0.4] {$Y(s)$} (out);

    %% Lazo de realimentación
    \draw[-]  (yout)  |- (junc);
    \draw[->] (junc)  -- (sensor);
    \draw[->] (sensor) -| node[pos=0.95, left] {$-$}
                          node[near end, below] {$B(s)$} (suma);

    %% Signo positivo del sumador
    \node[above left=0.15cm of suma, font=\small] {$+$};
  \end{tikzpicture}
  \caption{Diagrama de bloques general de un sistema de control de
           \textbf{lazo cerrado con realimentación negativa}.
           El signo $-$ en el sumador indica que la señal de realimentación
           se resta de la referencia para generar el error.}
  \label{fig:lazo-cerrado}
\end{figure}

\subsection{Mecanismo de la realimentación negativa}

El mecanismo de corrección opera de la siguiente manera, paso a paso:

\begin{enumerate}
  \item La salida $Y(s)$ es medida por el sensor, generando $B(s) = H(s) \cdot Y(s)$.
  \item El sumador calcula $E(s) = R(s) - B(s)$.
  \item Si $Y(s) > R(s)/H(s)$ (salida demasiado alta), entonces $E(s) < 0$ y el
    controlador reduce su acción $U(s)$, frenando la planta.
  \item Si $Y(s) < R(s)/H(s)$ (salida demasiado baja), entonces $E(s) > 0$ y el
    controlador aumenta su acción $U(s)$, empujando la planta hacia el valor deseado.
  \item El proceso se repite continuamente, manteniendo el error pequeño.
\end{enumerate}

Este mecanismo es \textbf{negativo} (se resta la realimentación) porque actúa
\textit{en contra} de la desviación: si la salida sube demasiado, la acción
correctiva la baja. Esto es lo que garantiza la estabilidad del lazo.

\interfisica{La realimentación negativa es el principio que subyace en
innumerables sistemas naturales y artificiales: el termostato del hogar,
la homeostasis del cuerpo humano (temperatura, glucosa en sangre), el
regulador de velocidad de un motor, el AVR de un generador.}

\subsection{La función de transferencia de lazo cerrado}

Si los bloques tienen funciones de transferencia $G_c(s)$, $G_p(s)$ y $H(s)$,
podemos obtener la función de transferencia del lazo cerrado. Definimos:
\begin{equation}
  G(s) = G_c(s)\,G_p(s)
  \qquad \text{(función de transferencia de lazo directo)}
  \label{eq:lazo-directo}
\end{equation}

De la figura, la relación entre la salida y la referencia es:
\begin{align}
  Y(s)   &= G(s)\,E(s) \label{eq:lc-1}\\
  E(s)   &= R(s) - H(s)\,Y(s) \label{eq:lc-2}
\end{align}

Sustituyendo \eqref{eq:lc-2} en \eqref{eq:lc-1}:
\begin{align}
  Y(s) &= G(s)\bigl[R(s) - H(s)\,Y(s)\bigr] \notag\\
  Y(s) + G(s)\,H(s)\,Y(s) &= G(s)\,R(s) \notag\\
  Y(s)\bigl[1 + G(s)\,H(s)\bigr] &= G(s)\,R(s)
\end{align}

Por lo tanto, la función de transferencia de lazo cerrado es:

\begin{resultado}[Función de transferencia de lazo cerrado]
  \begin{equation}
    \frac{Y(s)}{R(s)} = \frac{G(s)}{1 + G(s)\,H(s)}
    = \frac{G_c(s)\,G_p(s)}{1 + G_c(s)\,G_p(s)\,H(s)}
    \label{eq:ftlc}
  \end{equation}
\end{resultado}

Esta es una de las fórmulas más importantes del curso. El denominador
$1 + G(s)H(s)$ se denomina \textbf{polinomio característico} del sistema,
y sus raíces (los \textbf{polos de lazo cerrado}) determinan la estabilidad
y la dinámica del sistema \verCap{6}.

Para el caso especial de \textbf{realimentación unitaria} ($H(s) = 1$):
\begin{equation}
  \frac{Y(s)}{R(s)} = \frac{G(s)}{1 + G(s)}
  \label{eq:ftlc-unitaria}
\end{equation}

\subsection{Ventajas del lazo cerrado}

\begin{itemize}
  \item \textbf{Rechazo de perturbaciones:} si $d(t)$ perturba la planta,
    el sensor lo detecta a través de la variación en $y(t)$ y el controlador
    genera la corrección. El sistema ``lucha'' contra la perturbación.
  \item \textbf{Reducción de la sensibilidad a variaciones de parámetros:}
    la realimentación reduce el efecto de las incertidumbres en el modelo
    de la planta (demostrado formalmente en capítulos posteriores).
  \item \textbf{Mejora de la linealidad:} la no linealidad de la planta
    se reduce con realimentación de alta ganancia.
  \item \textbf{Posibilidad de modificar la dinámica:} el controlador permite
    acelerar la respuesta, aumentar el amortiguamiento o eliminar el error
    estacionario, aunque la planta por sí misma sea lenta o tenga error.
\end{itemize}

\subsection{Desventajas del lazo cerrado}

\begin{itemize}
  \item \textbf{Posibilidad de inestabilidad:} si el controlador no está
    bien diseñado, la realimentación puede provocar oscilaciones que crecen
    hasta saturar los actuadores. Un sistema inestable es inutilizable y,
    en la práctica, peligroso.
  \item \textbf{Mayor complejidad y costo:} se necesitan sensores, comparadores,
    cableado adicional y algoritmos de control.
  \item \textbf{Ruido de medición amplificado:} el ruido del sensor ingresa
    al lazo y puede contaminar la acción de control, especialmente si se
    usa acción derivativa.
\end{itemize}

%======================================================================
\section{Comparación: lazo abierto vs.\ lazo cerrado}
\label{sec:intro-comparacion}
%======================================================================

La \cref{tab:oa-vs-oc} resume las diferencias entre ambos tipos de sistemas.
La elección entre uno y otro depende del nivel de precisión requerido, el
costo admisible y las características del proceso.

\begin{table}[htbp]
  \centering
  \caption{Comparación entre sistemas de control de lazo abierto y lazo cerrado.}
  \label{tab:oa-vs-oc}
  \begin{tabularx}{\textwidth}{@{} l X X @{}}
    \toprule
    \textbf{Característica} & \textbf{Lazo abierto} & \textbf{Lazo cerrado} \\
    \midrule
    Necesidad de sensor & No & Sí (necesario) \\
    Costo y complejidad & Bajo & Mayor \\
    Precisión ante perturbaciones & Baja & Alta \\
    Sensibilidad a parámetros & Alta & Reducida \\
    Error en estado estacionario & Generalmente alto & Puede ser nulo (con integrador) \\
    Estabilidad & Siempre (si la planta es estable) & Puede ser inestable si mal diseñado \\
    Velocidad de respuesta & Depende solo de la planta & Puede mejorarse con el controlador \\
    Sensibilidad al ruido de medición & No afectado & Puede afectar \\
    Aplicaciones típicas & Temporizadores, semáforos fijos,
                           lavadoras sin sensor & Servo-motores, reguladores de tensión,
                           control de temperatura, AVR \\
    \bottomrule
  \end{tabularx}
\end{table}

\begin{observacion}
  En la práctica industrial moderna, la gran mayoría de los sistemas de
  control son de lazo cerrado. El lazo abierto se reserva para procesos
  donde la precisión no es crítica, la perturbación es muy pequeña o el
  sensor es demasiado costoso. Un buen ingeniero debe saber identificar
  cuándo cada opción es la apropiada.
\end{observacion}

%======================================================================
\section{Especificaciones de desempeño}
\label{sec:intro-especificaciones}
%======================================================================

Cuando se diseña un sistema de control, no basta con que la salida
``eventualmente'' llegue al valor deseado. Se especifican criterios cuantitativos
que el sistema debe satisfacer. Estas especificaciones se dividen en tres grupos:

\subsection{Estabilidad}

La estabilidad es la especificación más fundamental. Un sistema es
\textbf{estable} (BIBO: \textit{Bounded Input Bounded Output}) si para
toda entrada acotada la salida es también acotada. En términos prácticos:
si perturbamos el sistema y luego lo dejamos libre, la salida debe volver
a un valor de equilibrio en lugar de crecer indefinidamente.

\begin{definicion}{Estabilidad BIBO}
  Un sistema lineal e invariante en el tiempo (LTI) es \textbf{estable} si
  y solo si todos los polos de su función de transferencia de lazo cerrado
  tienen parte real negativa (están en el semiplano izquierdo del plano
  complejo $s$).
\end{definicion}

Un sistema inestable no puede funcionar en la práctica: la salida crece hasta
saturar los actuadores, lo que generalmente destruye el equipo. \textbf{La
estabilidad es una condición necesaria antes de evaluar cualquier otra
especificación.}

\subsection{Especificaciones en el dominio del tiempo}

Una vez garantizada la estabilidad, se evalúa la \textbf{respuesta transitoria}:
cómo llega el sistema al valor deseado.

% [INSERTAR FIGURA: Respuesta al escalón típica de un sistema de segundo orden
%  subamortiguado. Marcar claramente: tiempo de subida tr, tiempo pico tp,
%  sobreoscilación Mp, tiempo de establecimiento ts, valor final y_final.
%  El alumno debe identificar estas magnitudes y entender que cada una
%  corresponde a una especificación de diseño diferente.]
\figplaceholder{Respuesta al escalón de un sistema de segundo orden.
Se muestran: tiempo de subida $t_r$, tiempo pico $t_p$, sobreoscilación
porcentual $M_p$, tiempo de establecimiento $t_s$ y valor final $y_\infty$.}

Las especificaciones temporales más importantes son:

\begin{description}
  \item[Tiempo de subida ($t_r$):] tiempo que tarda la respuesta en pasar
    del 10\% al 90\% de su valor final. Cuantifica la \textit{rapidez} de
    la respuesta. Un $t_r$ pequeño significa un sistema rápido.
  \item[Tiempo pico ($t_p$):] tiempo en que la respuesta alcanza su primer
    máximo (sobreoscilación). Solo existe en sistemas subamortiguados.
  \item[Sobreoscilación porcentual ($M_p$):] porcentaje que la respuesta
    supera el valor final antes de establecerse:
    \begin{equation}
      M_p = \frac{y(t_p) - y(\infty)}{y(\infty)} \times 100\%
    \end{equation}
    Un $M_p$ alto indica escaso amortiguamiento. En muchas aplicaciones
    se especifica $M_p \leq 5\%$ o $M_p \leq 10\%$.
  \item[Tiempo de establecimiento ($t_s$):] tiempo a partir del cual la
    respuesta permanece dentro de una banda de $\pm 2\%$ (o $\pm 5\%$)
    alrededor del valor final. Cuantifica cuándo el transitorio ha terminado.
  \item[Error en estado estacionario ($e_{ss}$):] diferencia entre la
    referencia y la salida una vez que el transitorio ha desaparecido:
    \begin{equation}
      e_{ss} = \lim_{t\to\infty} e(t) = \lim_{t\to\infty} \bigl[r(t) - y(t)\bigr]
    \end{equation}
    Un controlador bien diseñado logra $e_{ss} = 0$ para entradas de
    referencia de tipo escalón o rampa \verCap{5}.
\end{description}

\subsection{Especificaciones en el dominio de la frecuencia}

Además del dominio temporal, se usan especificaciones en el dominio de la
frecuencia, relacionadas con la respuesta del sistema a señales sinusoidales:

\begin{description}
  \item[Margen de ganancia (MG):] cuánto se puede aumentar la ganancia
    del lazo antes de que el sistema se vuelva inestable. Se mide en decibeles.
    Un valor típico requerido es $MG \geq \SI{6}{\decibel}$.
  \item[Margen de fase (MF):] cuánto retardo de fase puede soportar el lazo
    antes de inestabilizarse. Un valor típico requerido es $MF \geq 30\grados$.
  \item[Ancho de banda ($\omega_{BW}$):] rango de frecuencias para el que
    el sistema responde adecuadamente. Relacionado con la velocidad de respuesta.
\end{description}

Estos conceptos se desarrollan en detalle en los capítulos de respuesta en
frecuencia \verCap{9} y diagramas de Bode \verCap{10}.

\begin{nota}[Conflicto entre especificaciones]
  Las especificaciones de desempeño están en conflicto entre sí.
  Un sistema más rápido (menor $t_r$) tiende a tener mayor sobreoscilación
  y menor margen de fase. El diseño de controladores es un problema de
  compromiso (\textit{trade-off}) entre especificaciones contradictorias.
\end{nota}

%======================================================================
\section{Clasificación de sistemas de control}
\label{sec:intro-clasificacion}
%======================================================================

Los sistemas de control pueden clasificarse desde distintos puntos de vista:

\subsection{Por el tipo de señal}

\begin{description}
  \item[Sistemas continuos (tiempo continuo):] las señales son funciones
    continuas del tiempo. Se analizan con la transformada de Laplace.
    Este es el enfoque principal del curso.
  \item[Sistemas discretos (tiempo discreto):] las señales son muestreadas
    a intervalos regulares. Se analizan con la transformada $\mathcal{Z}$.
    Son propios de los sistemas de control digital implementados en microcontroladores.
\end{description}

\subsection{Por el tipo de modelo}

\begin{description}
  \item[Sistemas lineales e invariantes en el tiempo (LTI):] satisfacen los
    principios de superposición y homogeneidad. Sus parámetros no cambian con
    el tiempo. Son el objeto de estudio del curso: ecuaciones diferenciales
    lineales de coeficientes constantes.
  \item[Sistemas no lineales:] no satisfacen la superposición. En la práctica,
    todos los sistemas reales son no lineales en algún grado. Para rangos de
    operación limitados, la linealización en torno a un punto de operación
    permite aplicar las herramientas del control lineal.
  \item[Sistemas variantes en el tiempo (LTV):] los parámetros del modelo
    cambian con el tiempo. Requieren técnicas de control adaptivo.
\end{description}

\subsection{Por la cantidad de entradas y salidas}

\begin{description}
  \item[SISO (\textit{Single Input Single Output}):] una sola variable de
    entrada controlable y una sola variable de salida medida. Es el caso
    que se estudia en la mayor parte del curso.
  \item[MIMO (\textit{Multiple Input Multiple Output}):] múltiples entradas
    y salidas, con posibles interacciones entre ellas. Requiere el enfoque
    de variables de estado \verCap{14}.
\end{description}

\subsection{Por el tipo de acción de control}

\begin{description}
  \item[Control de dos posiciones (on/off):] el actuador solo tiene dos
    estados: encendido o apagado. Ejemplo: termostato doméstico, interruptor
    automático. Simple pero introduce oscilaciones permanentes.
  \item[Control proporcional (P):] la acción de control es proporcional
    al error. Reduce el error pero no lo elimina completamente.
  \item[Control integral (I):] integra el error en el tiempo. Elimina el
    error estacionario pero puede reducir la estabilidad.
  \item[Control PID:] combinación de los tres anteriores. Es el estándar
    industrial \verCap{13}.
\end{description}

%======================================================================
\section{Ejemplo integrador: control de velocidad de un motor DC}
\label{sec:intro-ejemplo-motor}
%======================================================================

Para ilustrar todos los conceptos anteriores, analizamos el control de
velocidad de un motor de corriente continua (DC), uno de los sistemas
de control más frecuentes en la ingeniería eléctrica.

\subsection{Descripción del sistema}

Un motor DC alimentado por tensión $v_a$ produce un par electromagnético
que hace girar una carga mecánica. La velocidad angular $\omega(t)$ depende
del equilibrio entre el par motor y el par resistente de la carga.

Las variables del sistema son:
\begin{itemize}
  \item \textbf{Variable controlada:} velocidad angular $\omega(t)$
    [\si{\radian\per\second}] o $n(t)$ [RPM].
  \item \textbf{Variable de referencia:} velocidad deseada $\omega^*(t)$.
  \item \textbf{Variable de control:} tensión de armadura $v_a(t)$ [\si{\volt}].
  \item \textbf{Perturbación:} par de carga $T_c(t)$ [\si{\newton\meter}].
\end{itemize}

\subsection{Modelo simplificado}

Para este ejemplo introductorio, usamos el modelo de primer orden del motor
(sin inductancia de armadura):
\begin{equation}
  \tau_m \frac{d\omega}{dt} + \omega = K_m\,v_a - K_c\,T_c
  \label{eq:motor-dc}
\end{equation}
donde $\tau_m$ es la constante de tiempo mecánica, $K_m$ la ganancia del
motor y $K_c$ la constante de perturbación de carga.

En el dominio de Laplace (con condiciones iniciales nulas):
\begin{equation}
  G_p(s) = \frac{\Omega(s)}{V_a(s)} = \frac{K_m}{\tau_m\,s + 1}
  \label{eq:ft-motor}
\end{equation}
Este es el modelo de un sistema de \textbf{primer orden}, el más simple posible.

\subsection{Análisis de lazo abierto}

En lazo abierto, la velocidad en estado estacionario ante una tensión
constante $V_a = V_0$ es:
\begin{equation}
  \omega_{ss}^{OL} = K_m\,V_0 - K_c\,T_c
  \label{eq:motor-oa}
\end{equation}

Si la carga varía ($T_c$ cambia), la velocidad cambia proporcionalmente.
No hay ningún mecanismo automático de corrección.

\subsection{Control de lazo cerrado con sensor tacométrico}

Si instalamos un tacómetro (sensor de velocidad con ganancia $K_t$
[\si{\volt\per\radian\per\second}]) y un controlador proporcional de
ganancia $K_c^p$, el lazo cerrado produce:

\begin{align}
  E(s) &= \Omega^*(s) - K_t\,\Omega(s) \\
  V_a(s) &= K_c^p\,E(s)
\end{align}

La función de transferencia de lazo cerrado resulta:
\begin{equation}
  \frac{\Omega(s)}{\Omega^*(s)} =
  \frac{K_c^p\,K_m/\tau_m}{s + (1 + K_c^p\,K_m\,K_t)/\tau_m}
  \label{eq:motor-lc}
\end{equation}

La constante de tiempo de lazo cerrado es:
\begin{equation}
  \tau_{CL} = \frac{\tau_m}{1 + K_c^p\,K_m\,K_t}
  \label{eq:tau-cl}
\end{equation}

\interfisica{La constante de tiempo de lazo cerrado es \textbf{menor} que
la de lazo abierto en un factor $(1 + K_c^p K_m K_t)$. Esto significa que
el lazo cerrado \textbf{responde más rápido} que la planta sola. Cuanto
mayor sea la ganancia del controlador $K_c^p$, más rápida es la respuesta.
Sin embargo, como veremos en capítulos posteriores, existe un límite práctico
a cuánto se puede aumentar la ganancia antes de comprometer la estabilidad.}

% [INSERTAR FIGURA: Diagrama de bloques del control de velocidad de motor DC
%  en lazo cerrado. Mostrar: referencia de velocidad omega*, sumador, controlador
%  proporcional Kcp, modelo del motor Km/(tau_m*s+1), salida omega, tacómetro Kt.
%  Marcar la perturbación de carga Tc. El alumno debe identificar todos los
%  componentes enumerados en la sección anterior.]
\figplaceholder{Diagrama de bloques del control de velocidad de motor DC
en lazo cerrado. Referencia $\omega^*$, controlador proporcional $K_c^p$,
planta $K_m/(\tau_m s + 1)$, tacómetro $K_t$, perturbación $T_c$.}

%======================================================================
\section{Ejercicios resueltos}
\label{sec:intro-ejercicios-resueltos}
%======================================================================

\begin{ejercicioresuelto}{Identificación de componentes de un sistema de control (Nivel básico)}

\textbf{Enunciado.} Un horno industrial de tratamiento térmico usa un controlador
electrónico para mantener la temperatura a \SI{800}{\celsius}. Un termopar
mide la temperatura real. El calentamiento se realiza mediante resistencias
eléctricas alimentadas por un tiristor que regula la potencia. Identificá
los componentes del sistema de control y trazá un diagrama de bloques.

\textbf{Análisis conceptual.} Debemos identificar: planta, actuador, sensor,
controlador, referencia, variable controlada y señal de error. La descripción
contiene todos los componentes típicos de un lazo cerrado.

\textbf{Desarrollo.}
\begin{itemize}
  \item \textbf{Variable controlada (salida):} temperatura del horno $T(t)$ [\si{\celsius}].
  \item \textbf{Planta:} el horno industrial (masa metálica, aislamiento,
    geometría interna). Convierte la potencia eléctrica en calor y la distribuye.
  \item \textbf{Actuador:} el circuito de tiristores (rectificador controlado
    de silicio, SCR). Convierte la señal de control de baja potencia (ángulo
    de disparo) en la potencia eléctrica de las resistencias.
  \item \textbf{Sensor (transductor):} el termopar tipo K (o tipo R para altas
    temperaturas). Convierte la temperatura en una señal de tensión milimétrica.
  \item \textbf{Controlador:} el controlador electrónico de temperatura (un
    PID industrial). Procesa el error y genera el ángulo de disparo del tiristor.
  \item \textbf{Señal de referencia:} la temperatura de consigna \SI{800}{\celsius},
    introducida por el operario.
  \item \textbf{Señal de error:} $e(t) = T^*(t) - T_{medida}(t)$.
  \item \textbf{Perturbaciones:} apertura de la puerta del horno, variaciones
    en el material a tratar, fluctuaciones de la tensión de red.
\end{itemize}

El diagrama de bloques es:
\begin{center}
\begin{tikzpicture}[auto, node distance=2.4cm, thick, font=\small]
  \node[input] (ref) {};
  \node[sumnode, right of=ref, node distance=1.7cm] (suma) {};
  \node[block, right of=suma, node distance=3.0cm] (ctrl) {Controlador\\PID};
  \node[block, right of=ctrl, node distance=3.2cm] (act) {Tiristor\\(Actuador)};
  \node[block, right of=act, node distance=3.2cm] (horno) {Horno\\(Planta)};
  \node[output, right of=horno, node distance=2.2cm] (out) {};
  \node[block, below of=act, node distance=1.7cm] (sensor) {Termopar\\(Sensor)};
  \node[junction, below of=out, node distance=1.7cm] (junc) {};

  \draw[->] (ref)   -- node[above, pos=0.5] {$T^*$} (suma);
  \draw[->] (suma)  -- node[above] {$e(t)$} (ctrl);
  \draw[->] (ctrl)  -- node[above, font=\tiny, align=center] {ángulo de\\disparo} (act);
  \draw[->] (act)   -- node[above, font=\tiny, align=center] {potencia\\{[kW]}} (horno);
  \draw[->] (horno) -- node[above, name=temp, pos=0.4] {$T(t)$} (out);
  \draw[-]  (temp)  |- (junc);
  \draw[->] (junc)  -- (sensor);
  \draw[->] (sensor) -| node[pos=0.94, left] {$-$}
                        node[near end, below, font=\tiny] {$T_{med}$} (suma);
  \node[above left=0.12cm of suma, font=\tiny] {$+$};
\end{tikzpicture}
\end{center}

\textbf{Interpretación.} Este sistema de lazo cerrado mantiene la temperatura
a \SI{800}{\celsius} incluso si alguien abre la puerta del horno (perturbación):
el termopar detecta la caída de temperatura, el error se hace positivo, el
controlador aumenta la potencia y el horno recupera la temperatura. En lazo
abierto, eso no sería posible.
\end{ejercicioresuelto}

%% ---------------------------------------------------------------
\begin{ejercicioresuelto}{Comparación lazo abierto vs.\ lazo cerrado en un calentador eléctrico (Nivel intermedio)}

\textbf{Enunciado.} Un calentador eléctrico tiene los siguientes datos:
\begin{itemize}
  \item Potencia de calentamiento: $P = \SI{1200}{\watt}$
  \item Masa de agua: $m = \SI{4}{\kilogram}$
  \item Calor específico del agua: $c_p = \SI{4180}{\joule\per\kilogram\per\celsius}$
  \item Temperatura inicial: $T_0 = \SI{20}{\celsius}$
  \item Temperatura deseada: $T^* = \SI{65}{\celsius}$
\end{itemize}

\textbf{a)} En lazo abierto con temporizador fijo de $t = \SI{615}{\second}$:
¿qué temperatura se alcanza? ¿cuál es el error?

\textbf{b)} Si alguien usa \SI{1}{\kilogram} de agua (deja $m = \SI{3}{\kilogram}$),
¿qué temperatura alcanza el sistema en lazo abierto con el mismo temporizador?

\textbf{c)} Con un termostato (lazo cerrado), describí cualitativamente qué ocurre
en el caso b). ¿Cuál es el error estacionario?

\textbf{Análisis conceptual.} La energía entregada por el calentador en un
tiempo $t$ es $Q = P \cdot t$. La temperatura final se obtiene del balance
de energía: $Q = m\,c_p\,\Delta T$.

\textbf{Desarrollo.}

\textbf{Parte a)} Energía entregada:
\begin{equation*}
  Q = P\,t = 1200 \times 615 = \SI{738000}{\joule} = \SI{738}{\kilo\joule}
\end{equation*}
Aumento de temperatura:
\begin{equation*}
  \Delta T = \frac{Q}{m\,c_p} = \frac{738000}{4 \times 4180} = \frac{738000}{16720}
  \approx \SI{44.1}{\celsius}
\end{equation*}
Temperatura final:
\begin{equation*}
  T_f^{(a)} = T_0 + \Delta T = 20 + 44.1 = \SI{64.1}{\celsius}
\end{equation*}
Error en lazo abierto (caso nominal):
\begin{equation*}
  e_{OL}^{(a)} = T^* - T_f^{(a)} = 65 - 64.1 \approx \SI{0.9}{\celsius}
\end{equation*}
Pequeño, porque el temporizador fue calibrado para este caso.

\textbf{Parte b)} Misma energía $Q = \SI{738}{\kilo\joule}$, pero $m = \SI{3}{\kilogram}$:
\begin{equation*}
  \Delta T' = \frac{738000}{3 \times 4180} = \frac{738000}{12540}
  \approx \SI{58.9}{\celsius}
\end{equation*}
\begin{equation*}
  T_f^{(b)} = 20 + 58.9 = \SI{78.9}{\celsius}
\end{equation*}
Error en lazo abierto con la perturbación:
\begin{equation*}
  e_{OL}^{(b)} = 65 - 78.9 = \SI{-13.9}{\celsius}
\end{equation*}
El agua se sobrecalienta casi \SI{14}{\celsius} por encima del objetivo.
Con menos agua, el calentamiento es más rápido y el temporizador fijo
no puede adaptar su tiempo.

\textbf{Parte c)} Con termostato (lazo cerrado): el termostato mide la
temperatura continuamente. Cuando $T < \SI{65}{\celsius}$, el calentador
está encendido; cuando $T \geq \SI{65}{\celsius}$, se apaga.
Independientemente de si hay 3 kg o 4 kg de agua, el sistema se detiene
exactamente cuando alcanza \SI{65}{\celsius}. El error estacionario es cero
(o mínimo, según la histéresis del termostato).

\textbf{Interpretación.} Este ejemplo ilustra con claridad la principal
ventaja del lazo cerrado: \textbf{rechazo de perturbaciones}. El lazo abierto
funciona bien solo para las condiciones para las que fue calibrado. Cualquier
desviación en los parámetros produce un error no corregible.
\end{ejercicioresuelto}

%% ---------------------------------------------------------------
\begin{ejercicioresuelto}{Control proporcional de velocidad de motor DC (Tipo examen)}

\textbf{Enunciado.} Un motor DC tiene las siguientes características:
\begin{itemize}
  \item Ganancia estática: $K_m = \SI{3}{RPM/V}$
  \item Constante de tiempo mecánica: $\tau_m = \SI{0.4}{\second}$
  \item Ganancia del sensor tacométrico: $K_t = \SI{0.02}{V/RPM}$
\end{itemize}
La velocidad de referencia es $n^* = \SI{1000}{RPM}$.
El controlador es proporcional con ganancia $K_c^p = 50$.

\textbf{a)} ¿Cuál es la tensión de referencia equivalente?

\textbf{b)} Calculá la velocidad en estado estacionario en lazo cerrado
(sin perturbación) y el error porcentual.

\textbf{c)} Si la perturbación de carga produce una reducción de
\SI{200}{RPM} en lazo abierto, ¿cuánto reduce esa misma perturbación
la velocidad en lazo cerrado?

\textbf{d)} ¿Cuál es la constante de tiempo del lazo cerrado?

\textbf{Desarrollo.}

\textbf{Parte a)} La referencia eléctrica equivalente:
\begin{equation*}
  V_{ref} = \frac{n^*}{K_m} = \frac{1000}{3} \approx \SI{333.3}{\volt}
\end{equation*}
(Esta es la tensión que, en lazo abierto, produciría la velocidad nominal.
En lazo cerrado, la referencia se expresa como la consigna de velocidad
directamente, y el bloque comparador hace la resta en RPM.)

\textbf{Parte b)} Ganancia de lazo abierto:
\begin{equation*}
  G_0 = K_c^p \cdot K_m \cdot K_t = 50 \times 3 \times 0.02 = 3
\end{equation*}
Velocidad en estado estacionario de lazo cerrado con referencia $n^* = \SI{1000}{RPM}$:
\begin{equation*}
  n_{ss}^{CL} = \frac{K_c^p \cdot K_m}{1 + K_c^p \cdot K_m \cdot K_t} \cdot
  \frac{n^*}{K_m/1}
\end{equation*}
Usando la fórmula de lazo cerrado con ganancia $G_0$:
\begin{equation*}
  n_{ss}^{CL} = \frac{G_0}{1 + G_0} \cdot n^* = \frac{3}{1+3} \cdot 1000
  = \SI{750}{RPM}
\end{equation*}
Error estacionario en lazo cerrado:
\begin{equation*}
  e_{ss}^{CL} = n^* - n_{ss}^{CL} = 1000 - 750 = \SI{250}{RPM}
  \quad \Rightarrow \quad 25\%
\end{equation*}

\textbf{Parte c)} Efecto de la perturbación. En lazo abierto, la perturbación
produce $\Delta n_{OL} = \SI{-200}{RPM}$.

En lazo cerrado, la perturbación se atenúa por el factor $\dfrac{1}{1 + G_0}$:
\begin{equation*}
  \Delta n_{CL} = \frac{\Delta n_{OL}}{1 + G_0} = \frac{-200}{1 + 3}
  = \SI{-50}{RPM}
\end{equation*}
La perturbación solo produce $\SI{-50}{RPM}$ en lazo cerrado frente a
$\SI{-200}{RPM}$ en lazo abierto: una mejora de \textbf{4 veces} (igual al
factor $1 + G_0 = 4$).

\textbf{Parte d)} Constante de tiempo del lazo cerrado:
\begin{equation*}
  \tau_{CL} = \frac{\tau_m}{1 + G_0} = \frac{0.4}{1 + 3} = \SI{0.1}{\second}
\end{equation*}
El lazo cerrado es 4 veces más rápido que la planta sola.

\textbf{Interpretación.} Los resultados demuestran las fortalezas y limitaciones
del control proporcional:
\begin{itemize}
  \item \textbf{Fortaleza:} mejora el rechazo de perturbaciones (factor $1 + G_0$)
    y acelera la respuesta.
  \item \textbf{Limitación:} no elimina el error estacionario. Con $K_c^p = 50$
    y $G_0 = 3$, el error es del 25\%. Para eliminarlo habría que usar acción
    integral (controlador PI o PID) \verCap{13}.
\end{itemize}
Aumentar $K_c^p$ reduce el error estacionario, pero tiene un límite: a partir de
cierta ganancia el sistema se vuelve oscilatorio e inestable. Encontrar el
balance entre precisión y estabilidad es el problema central del diseño de
controladores.
\end{ejercicioresuelto}

%======================================================================
\section{Ejercicios propuestos}
\label{sec:intro-propuestos}
%======================================================================

\begin{enumerate}
  \item Un sistema de control de nivel de agua en un tanque industrial usa un
    flotante (sensor), una válvula motorizada (actuador) y un controlador PLC.
    El nivel deseado es de \SI{2.5}{\meter}.
    \begin{enumerate}
      \item Identificá todos los componentes del sistema.
      \item Trazá el diagrama de bloques indicando señales y bloques.
      \item ¿Qué tipo de perturbaciones puede afectar el sistema?
      \item ¿Qué ocurriría si el flotante falla (queda en la posición de ``nivel
        correcto'')? ¿Qué tipo de falla provoca en el lazo de control?
    \end{enumerate}

  \item Un regulador automático de tensión (AVR) de un generador sincrónico
    de \SI{10}{\mega\volt\ampere} mantiene la tensión terminal en
    \SI{13.8}{\kilo\volt}. El transformador de medida de tensión (TV) tiene
    una relación de \num{13800/110}.
    \begin{enumerate}
      \item ¿Cuál es la tensión de referencia en el secundario del TV?
      \item Si el TV tiene un error del 0.5\%, ¿cuál es el error máximo en
        la tensión terminal del generador?
      \item Proponé un esquema de lazo cerrado para este sistema.
    \end{enumerate}

  \item Para el motor DC del Ejercicio Resuelto 1.3, si se desea un error
    estacionario máximo del 2\%:
    \begin{enumerate}
      \item ¿Cuál es la ganancia de lazo mínima $G_0$ necesaria?
      \item ¿Cuál debe ser $K_c^p$ para lograr esa ganancia?
      \item ¿Cuánto se reduce la perturbación de lazo con esa ganancia?
    \end{enumerate}
    \textit{Pista: $e_{ss}\% = \frac{1}{1+G_0} \times 100\%$.}

  \item Un inversor fotovoltaico debe inyectar corriente senoidal a la red
    (220 V, 50 Hz). El control de corriente tiene una referencia de corriente
    $i^*(t) = I_m\,\sen(\omega t)$. Describí cualitativamente:
    \begin{enumerate}
      \item ¿Por qué es necesario un lazo cerrado (y no lazo abierto)?
      \item ¿Qué sensor se usaría para medir la corriente inyectada?
      \item ¿Qué tipo de perturbaciones afectan este sistema?
    \end{enumerate}

  \item \textbf{(Nivel avanzado)} Un sistema de control tiene la función
    de transferencia de lazo cerrado:
    \begin{equation*}
      T(s) = \frac{10}{s^2 + 4s + 10}
    \end{equation*}
    \begin{enumerate}
      \item Sin calcular nada, indicá si el sistema es estable (analizá los
        signos de los coeficientes del denominador).
      \item ¿Es este un sistema de primer o segundo orden? ¿Cómo lo sabés?
      \item ¿Qué ganancia estática tiene este sistema? \textit{Pista: evaluá $T(0)$.}
      \item ¿Cuál es el error estacionario ante una entrada escalón unitario?
    \end{enumerate}
\end{enumerate}

%======================================================================
\section{Implementación en MATLAB}
\label{sec:intro-matlab}
%======================================================================

Esta sección introduce las herramientas básicas de MATLAB para análisis de
sistemas de control. Se asume que se tiene instalado MATLAB con el
\textbf{Control System Toolbox}. Todo el código presentado está disponible
en el archivo \texttt{matlab/cap01\_introduccion.m}.

\subsection{Creación de funciones de transferencia}

El objeto fundamental en MATLAB para sistemas LTI es la función de
transferencia, creada con el comando \texttt{tf()}:

\begin{matlabcodetitulo}{Creación de funciones de transferencia en MATLAB}
% =========================================================
% Archivo: matlab/cap01_introduccion.m
% Objetivo: Introducción básica al Control System Toolbox
% Capítulo 1: Introducción a los sistemas de control
% =========================================================

%% 1. Definición de funciones de transferencia
% Sintaxis: G = tf(numerador, denominador)
% Los vectores contienen los coeficientes de mayor a menor grado.

% Motor DC: G(s) = Km / (tau_m * s + 1)
Km   = 3;      % ganancia estatica [RPM/V]
tau_m = 0.4;   % constante de tiempo [s]

G_motor = tf(Km, [tau_m, 1]);

disp('Funcion de transferencia del motor DC (lazo abierto):')
G_motor

% Sistema de segundo orden: G(s) = 10 / (s^2 + 4s + 10)
G_2orden = tf(10, [1, 4, 10]);

disp('Funcion de transferencia de segundo orden:')
G_2orden
\end{matlabcodetitulo}

El resultado que muestra MATLAB es:
\begin{verbatim}
Funcion de transferencia del motor DC (lazo abierto):

       3
  -----------
  0.4 s + 1
\end{verbatim}

\subsection{Respuesta al escalón}

El comando \texttt{step()} grafica la respuesta del sistema ante una entrada
escalón unitario ($r(t) = u(t)$):

\begin{matlabcodetitulo}{Respuesta al escalón: lazo abierto vs. lazo cerrado}
%% 2. Respuesta al escalon: lazo abierto vs. cerrado

Kcp = 50;      % ganancia del controlador proporcional
Kt  = 0.02;    % ganancia del tacometro [V/RPM]

% Controlador proporcional
G_ctrl = tf(Kcp, 1);

% Sensor (ganancia constante)
H_sensor = tf(Kt, 1);

% Lazo abierto (solo la planta)
G_LA = G_ctrl * G_motor;

% Lazo cerrado con realimentacion negativa
% feedback(G, H) calcula G/(1 + G*H)
G_LC = feedback(G_LA, H_sensor);

disp('Funcion de transferencia de lazo cerrado:')
G_LC

% Tiempo de simulacion
t = 0:0.001:1.5;    % 0 a 1.5 segundos, paso 1 ms

% Referencia: escalon de 1000 RPM
n_ref = 1000;

%% Grafica comparativa
figure('Name', 'Cap. 1 - Lazo Abierto vs Lazo Cerrado', ...
       'NumberTitle', 'off', 'Color', 'white');

% Lazo abierto (normalizado para ver la forma)
[y_LA, t_LA] = step(n_ref * G_motor / dcgain(G_motor), t);
subplot(1, 2, 1)
plot(t_LA, y_LA, 'b-', 'LineWidth', 2)
hold on
yline(n_ref, 'r--', 'Referencia', 'LineWidth', 1.5)
title('Lazo Abierto (normalizado)', 'FontSize', 12)
xlabel('Tiempo (s)')
ylabel('Velocidad (RPM)')
grid on; ylim([0, 1200])

% Lazo cerrado
[y_LC, t_LC] = step(n_ref * G_LC / dcgain(G_LC), t);
subplot(1, 2, 2)
plot(t_LC, y_LC, 'b-', 'LineWidth', 2)
hold on
yline(n_ref, 'r--', 'Referencia', 'LineWidth', 1.5)
title('Lazo Cerrado (K_c = 50)', 'FontSize', 12)
xlabel('Tiempo (s)')
ylabel('Velocidad (RPM)')
grid on; ylim([0, 1200])

sgtitle('Respuesta al escalon: Lazo Abierto vs. Lazo Cerrado', 'FontSize', 14)
\end{matlabcodetitulo}

\textbf{Qué se espera observar:}
\begin{itemize}
  \item Ambas curvas son exponenciales crecientes (sistema de primer orden).
  \item La curva de lazo cerrado alcanza el 63\% de su valor final en $\tau_{CL} = \SI{0.1}{\second}$,
    4 veces más rápida que la de lazo abierto ($\tau_m = \SI{0.4}{\second}$).
  \item El valor final de lazo cerrado ($n_{ss}^{CL} = \SI{750}{RPM}$) es menor
    que la referencia ($\SI{1000}{RPM}$): este es el error estacionario del control proporcional.
  \item El lazo abierto (normalizado) llega al 100\% de su valor final.
\end{itemize}

\subsection{Información de la respuesta al escalón}

MATLAB puede calcular automáticamente los parámetros de la respuesta:

\begin{matlabcodetitulo}{Análisis de la respuesta al escalón}
%% 3. Informacion cuantitativa de la respuesta

% stepinfo() devuelve una estructura con todos los parametros
info_LA = stepinfo(G_motor);
info_LC = stepinfo(G_LC);

fprintf('=== LAZO ABIERTO ===\n')
fprintf('Constante de tiempo: tau = %.3f s\n', tau_m)
fprintf('Tiempo de subida   : tr  = %.3f s\n', info_LA.RiseTime)
fprintf('Tiempo establecimiento: ts = %.3f s\n\n', info_LA.SettlingTime)

fprintf('=== LAZO CERRADO (Kcp = %d) ===\n', Kcp)
fprintf('Ganancia de lazo   : G0  = %.1f\n', Kcp * Km * Kt)
fprintf('Constante de tiempo: tau_CL = %.3f s\n', tau_m / (1 + Kcp*Km*Kt))
fprintf('Tiempo de subida   : tr  = %.3f s\n', info_LC.RiseTime)
fprintf('Tiempo establecimiento: ts = %.3f s\n', info_LC.SettlingTime)
fprintf('Ganancia CC (DC gain): %.4f\n', dcgain(G_LC))
fprintf('Error estacionario: %.1f %%\n', (1 - dcgain(G_LC)) * 100)
\end{matlabcodetitulo}

\begin{observacion}[Nota sobre normalización en MATLAB]
  En el código anterior se normaliza la respuesta al escalón del lazo
  abierto para comparar la \textit{forma} de la curva, no su valor final.
  En la práctica, el lazo abierto con $V_a = \SI{333}{\volt}$ produciría
  \SI{1000}{RPM} (exactamente la referencia), pero solo bajo las condiciones
  nominales. Sin perturbaciones, el lazo abierto da el valor correcto;
  con perturbaciones, no.
\end{observacion}

\subsection{Exploración: efecto de la ganancia del controlador}

El siguiente código permite visualizar cómo cambia la respuesta al variar
la ganancia proporcional $K_c^p$:

\begin{matlabcodetitulo}{Efecto de la ganancia proporcional en la respuesta}
%% 4. Efecto de la ganancia del controlador

figure('Name', 'Efecto de la ganancia proporcional', ...
       'Color', 'white');

Kcp_vals = [10, 25, 50, 100, 200];
colores = lines(length(Kcp_vals));
leyendas = {};

for i = 1:length(Kcp_vals)
    Kcp_i = Kcp_vals(i);
    G_LC_i = feedback(tf(Kcp_i * Km, [tau_m, 1]), H_sensor);
    G0_i   = Kcp_i * Km * Kt;
    ess_i  = 100 / (1 + G0_i);    % error estacionario [%]

    [y_i, t_i] = step(G_LC_i, t);
    y_i_rpm = y_i * 1000 / dcgain(G_LC_i);  % escalar a RPM reales

    plot(t_i, y_i_rpm, 'Color', colores(i,:), 'LineWidth', 1.8)
    hold on
    leyendas{end+1} = sprintf('K_c=%d  (e_{ss}=%.1f%%)', Kcp_i, ess_i);
end

yline(1000, 'k--', 'Referencia 1000 RPM', 'LineWidth', 1.5)
xlabel('Tiempo (s)', 'FontSize', 12)
ylabel('Velocidad (RPM)', 'FontSize', 12)
title('Efecto de K_c^p sobre la respuesta al escalon', 'FontSize', 13)
legend(leyendas, 'Location', 'southeast', 'FontSize', 9)
grid on
\end{matlabcodetitulo}

\textbf{Qué se espera observar:}
\begin{itemize}
  \item A mayor ganancia $K_c^p$, el sistema responde más rápido (menor $\tau_{CL}$).
  \item A mayor ganancia, el error estacionario disminuye: con $K_c^p = 200$
    el error es $\approx 1.4\%$.
  \item Todas las curvas son monótonas crecientes (sin sobreoscilación) porque
    el sistema es de primer orden, sin importar la ganancia.
  \item En un sistema de orden superior, al aumentar la ganancia la respuesta
    eventualmente se volvería oscilatoria e inestable \verCap{6}.
\end{itemize}

%======================================================================
\section{Gráficos sugeridos}
\label{sec:intro-graficos}
%======================================================================

% [INSERTAR FIGURA: Respuesta al escalón del motor DC en lazo abierto y
%  lazo cerrado (Kcp=50). Mostrar ambas curvas en el mismo gráfico para
%  comparación directa. Marcar tau_m=0.4s y tau_CL=0.1s. Mostrar el
%  error estacionario del lazo cerrado (750 RPM vs 1000 RPM).
%  Generado por: matlab/cap01_introduccion.m, sección 2.]
\figplaceholder{Respuesta al escalón: lazo abierto vs.\ lazo cerrado para
el motor DC. Eje $x$: tiempo [\si{\second}], eje $y$: velocidad [RPM].
Marcar $\tau_m = \SI{0.4}{\second}$ y $\tau_{CL} = \SI{0.1}{\second}$.
Error estacionario visible: $n_{ss}^{CL} = \SI{750}{RPM}$ vs $n^* = \SI{1000}{RPM}$.}

% [INSERTAR FIGURA: Efecto de la ganancia proporcional sobre la respuesta
%  al escalón. Mostrar 5 curvas para Kcp = 10, 25, 50, 100, 200.
%  Anotar el error estacionario de cada curva. El alumno observará la
%  disminución del tiempo de establecimiento y del error a medida que
%  aumenta la ganancia, sin que el sistema se vuelva inestable (1er orden).
%  Generado por: matlab/cap01_introduccion.m, sección 4.]
\figplaceholder{Efecto de la ganancia proporcional $K_c^p$ sobre la respuesta al
escalón del motor DC. Cinco curvas: $K_c^p \in \{10, 25, 50, 100, 200\}$.
Notar la reducción del error estacionario y el aumento de la velocidad de respuesta.}

% [INSERTAR FIGURA: Diagrama conceptual comparando 4 ejemplos de sistemas
%  de control en ingeniería eléctrica: (1) control de velocidad de motor,
%  (2) AVR de generador, (3) MPPT en solar, (4) control de corriente en inversor.
%  Para cada uno: mostrar planta, actuador, sensor y variable controlada.
%  Ilustración esquemática, no necesariamente diagrama de bloques formal.]
\figplaceholder{Cuatro ejemplos de sistemas de control en ingeniería eléctrica:
motor DC, generador con AVR, panel solar con MPPT, inversor con control de corriente.
Resaltar planta, sensor, actuador y variable controlada en cada caso.}

%======================================================================
\section{Resumen del capítulo}
\label{sec:intro-resumen}
%======================================================================

En este capítulo se establecieron los cimientos del control automático.
Los conceptos presentados son fundamentales para todo el resto del curso:

\begin{itemize}
  \item El \textbf{control automático} permite mantener una variable de un
    proceso en un valor deseado sin intervención humana continua, incluso
    frente a perturbaciones.
  \item El concepto central es la \textbf{realimentación negativa}: medir
    la salida, compararla con la referencia y actuar en función del error.
  \item Un sistema de \textbf{lazo cerrado} es superior al de lazo abierto
    en precisión y rechazo de perturbaciones, al costo de mayor complejidad
    y el riesgo de inestabilidad.
  \item Los componentes de un sistema de control son: \textbf{planta},
    \textbf{actuador}, \textbf{sensor}, \textbf{controlador} y \textbf{comparador}.
  \item Las \textbf{especificaciones de desempeño} se dividen en estabilidad,
    especificaciones transitorias ($t_r$, $t_p$, $M_p$, $t_s$) y error
    en régimen permanente ($e_{ss}$).
  \item La función de transferencia de lazo cerrado es
    $T(s) = G(s)/\bigl[1 + G(s)H(s)\bigr]$, y el denominador $1 + G(s)H(s)$
    determina la estabilidad del sistema.
  \item El \textbf{control proporcional} mejora la respuesta pero no elimina
    el error estacionario; para ello se necesita acción integral (PI o PID).
  \item \textbf{MATLAB} con el Control System Toolbox es la herramienta
    estándar para análisis y diseño de sistemas de control.
\end{itemize}

%======================================================================
\section{Conceptos clave}
\label{sec:intro-conceptos}
%======================================================================

\begin{conceptosclave}
\begin{description}[leftmargin=5.5em, labelwidth=5em, font=\bfseries\color{colkey}]
  \item[Planta]        Objeto físico cuya variable de salida se desea controlar.
  \item[Actuador]      Dispositivo que ejecuta la acción de control sobre la planta.
  \item[Sensor]        Transductor que mide la variable de salida y la convierte en señal eléctrica.
  \item[Error] $e(t) = r(t) - b(t)$. Diferencia entre la referencia y la medición de la salida.
  \item[Realim.] Uso de la señal de salida medida para generar la acción de control.
  \item[Lazo abierto]  Sistema donde la acción de control no depende de la salida.
  \item[Lazo cerrado]  Sistema donde la acción de control se genera a partir del error.
  \item[FT lazo cerr.] $T(s) = G(s)/[1+G(s)H(s)]$ con $G(s) = G_c(s)G_p(s)$.
  \item[$e_{ss}$]      Error en régimen permanente: $\lim_{t\to\infty}e(t)$.
  \item[Estabilidad]   Todos los polos de $T(s)$ con parte real negativa.
  \item[BIBO]          \textit{Bounded Input Bounded Output}: entrada acotada $\Rightarrow$ salida acotada.
  \item[$M_p$]         Sobreoscilación porcentual de la respuesta al escalón.
  \item[$t_s$]         Tiempo de establecimiento: cuando la respuesta queda dentro del $\pm 2\%$.
  \item[PID]           Controlador Proporcional-Integral-Derivativo. El más usado en industria.
\end{description}
\end{conceptosclave}

%======================================================================
\section{Errores comunes}
\label{sec:intro-errores}
%======================================================================

\begin{errorcomon}[Confundir ``lazo abierto'' con ``sin control'']
  Un sistema de lazo abierto \textbf{sí tiene controlador}; lo que no tiene
  es medición de la salida. Un microondas con temporizador es de lazo abierto:
  tiene un ``controlador'' (el temporizador) pero no mide la temperatura.
  ``Lazo abierto'' no significa ``sin control'', sino ``sin realimentación''.
\end{errorcomon}

\begin{errorcomon}[Creer que mayor ganancia siempre es mejor]
  Aumentar la ganancia del controlador proporcional reduce el error estacionario
  y acelera la respuesta, pero \textbf{no se puede aumentar indefinidamente}.
  En sistemas de orden 2 o superior, una ganancia excesiva provoca
  sobreoscilaciones crecientes y eventualmente inestabilidad. Siempre hay
  un compromiso entre precisión, velocidad y estabilidad.
\end{errorcomon}

\begin{errorcomon}[Omitir el signo negativo en la realimentación]
  En la suma del lazo cerrado, la señal realimentada se \textbf{resta}
  (signo negativo). Si se suma (signo positivo), la realimentación es
  \textit{positiva}, lo que generalmente inestabiliza el sistema. Este es
  un error frecuente al trazar diagramas de bloques en los exámenes.
\end{errorcomon}

\begin{errorcomon}[Confundir función de transferencia de lazo abierto y cerrado]
  La función de transferencia de lazo abierto es $G(s)H(s)$ (se abre el lazo
  y se mide la respuesta de ida y vuelta). La de lazo cerrado es
  $G(s)/[1+G(s)H(s)]$. Son distintas y se usan para propósitos diferentes:
  la de lazo abierto para análisis de estabilidad (criterio de Bode, Nyquist);
  la de lazo cerrado para calcular la respuesta a la referencia.
\end{errorcomon}

\begin{errorcomon}[Ignorar las perturbaciones al especificar el sistema]
  El error estacionario calculado a partir de la función de transferencia
  de referencia a salida solo considera la respuesta a $r(t)$. En un sistema
  real, las perturbaciones $d(t)$ también contribuyen al error.
  Un sistema de lazo cerrado reduce (pero no elimina con control P) el error
  causado por perturbaciones. Para análisis completo, se debe calcular
  también la función de transferencia de perturbación a salida.
\end{errorcomon}
